\subsection{テイラー級数}

有限テイラー公式は,剰余項を含めた等式である.これに対して
\[
 f(x)=\sum_{k=0}^{\infty}\frac{f^{(k)}(a)}{k!}(x-a)^k
\]
と書くには,\,$n\to\infty$のとき$R_n(x)\to0$になることを確かめる必要がある.有限の段階で正しい公式を何度でも書けることと,その極限で剰余項が消えることは別の主張である.

この違いは,等比数列の和にすでに現れている.\,$x\ne1$に対して
\[
 \frac1{1-x}=1+x+\cdots+x^n+\frac{x^{n+1}}{1-x}
\]
は有限の恒等式である.\,$|x|<1$なら最後の項は$0$に近づくので,無限級数の和も$\frac1{1-x}$になる.しかし$x=2$なら,最後の項は$-2^{n+1}$であり,小さくならない.有限の恒等式は正しいままであっても,剰余項を捨てることはできない.

テイラー公式で$|f^{(n+1)}(t)|\leq M_n$と評価したときも,同じことを調べる.\,$|x-a|^{n+1}$だけを見るのではなく,
\[
 \frac{M_n|x-a|^{n+1}}{(n+1)!}
\]
全体が$0$へ近づくかを調べなければならない.微分を繰り返すにつれて$M_n$が大きくなる場合もあるからである.

例えば$0\leq x\leq1$での$\log(1+x)$については,\,$(n+1)$階微分の絶対値が$n!$以下なので
\[
 |R_n(x)|\leq\frac{x^{n+1}}{n+1}\leq\frac1{n+1}.
\]
よって,この区間では剰余項が$0$に近づき,\,$x=1$も含めてテイラー級数が元の関数を表すことを証明できる.

\begin{note}
何回でも微分できる関数が,必ずテイラー級数と一致するわけでもない.例えば
\[
 h(x)=
 \begin{cases}
  \e^{-\frac1{x^2}}&(x\ne0),\\
  0&(x=0)
 \end{cases}
\]
は,\,$x=0$でのすべての階の微分係数が$0$になる.それでも$x\ne0$では$h(x)>0$である.そのため,\,$0$を中心に作ったテイラー多項式はすべて$0$であり,元の関数と一致しない.

この性質の理由は,指数関数の減少が任意のべきより速いことにある.\,$z>0$では,有限テイラー公式の剰余項の正値性から$\e^z\geq\frac{z^N}{N!}$である.\,$N$を十分大きく選べば,任意の固定した$m$について
\[
 \lim_{x\to0}|x|^{-m}\e^{-\frac1{x^2}}=0
\]
が従う.\,$x\ne0$で$h$を繰り返し微分すると,\,$x$の負の整数べきの有限和に$\e^{-\frac1{x^2}}$を掛けた形になる.この極限を,各階の微分係数の定義に順に使うと,\,$x=0$でどの階の微分係数も$0$になることが分かる.無限に滑らかであることと,テイラー級数で表せることは区別されるのである.
\end{note}



